\documentclass[preview]{standalone}

\usepackage{amsmath}
\usepackage{amssymb}
\usepackage{stellar}
\usepackage{enumitem}
\usepackage{definitions}

\begin{document}

\id{linearalgebra-exercises-batch-6}
\genpage

\section{Exercises}

\begin{snippetexercise}{linear-algebra-batch-6-ex-1}{}
    Determine whether the following collections of points in \(\mathbb{A}^3(\realnumbers)\) are affinely independent:\\
    \begin{align*}
        &\text{a) } \{(2, 1, 4), (4, 4, -1), (6, 7, -6)\}; &&\text{c) } \{(1, 2, 3), (5, -4, 7), (-1, -1, 6), (9, 2, 1)\};\\
        &\text{b) } \{(1, 2, 3), (-4, 2, 1), (1, 1, 2)\}; &&\text{d) } \{(3, 4, 5), (4, 5, 6), (1, 6, 6), (-1, 5, 8)\}.
    \end{align*}
\end{snippetexercise}

\begin{snippetsolution}{linear-algebra-batch-6-ex-1-sol}{}
    \todo
\end{snippetsolution}

\begin{snippetexercise}{linear-algebra-batch-6-ex-2}{}
    In the affine space \(\mathbb{A}^3(\realnumbers)\) consider the points \[P = (1, 1, 1), \quad Q = (3, 3, 5), \quad R = (-1, 0, 2) \quad \text{and} \quad S = (-3, 1, -2).\]
    \begin{enumerate}[label=\Alph*]
        \item Verify that \(\{P, Q, R, S\}\) is an affine frame of reference in \(\mathbb{A}^3(\realnumbers)\).
        \item Calculate the coordinates of \(O = (0, 0, 0)\) in this affine frame of reference.
    \end{enumerate}
\end{snippetexercise}

\begin{snippetsolution}{linear-algebra-batch-6-ex-2-sol}{}
    \todo
\end{snippetsolution}

\begin{snippetexercise}{linear-algebra-batch-6-ex-3}{}
    In each of the following cases determine, if it exists, the value of the parameter \(m \in \realnumbers\) for which a triplet of \emph{collinear} points in \(\mathbb{A}^3(\realnumbers)\) is obtained, that is, points lying on a line:\\
    \begin{align*}
        &\text{a) } \{(2, -1, 2), (1, 1, 1), (2, 1 - m, 4)\}; &&\text{c) } \{(1, -m, 0), (m, 1, 1), (1, -1, 3)\};\\
        &\text{b) } \{(3, 0, 0), (0, 1, 1), (m, m, m)\}; &&\text{d) } \{(1, m, 0), (2, \sqrt{2}, 1), (2, 110, 1)\}.
    \end{align*}
\end{snippetexercise}

\begin{snippetsolution}{linear-algebra-batch-6-ex-3-sol}{}
    \todo
\end{snippetsolution}

\begin{snippetexercise}{linear-algebra-batch-6-ex-4}{}
    In \(\mathbb{A}^4(\realnumbers)\) we fix the affine frame of reference \(\{0, e_1, e_2, e_3, e_4\}\) and consider the affine subspaces \[\pi = \{(x, y, z, t) \suchthat 2x + y - z + 3t = 2\} \quad \text{and} \quad \pi' = (1, 0, 0, 0) + \text{span}((3, 0, 1, 3), (2, -1, 0, -4)).\]
    Find parametric and Cartesian equations for \(\pi \cap \pi'\).
\end{snippetexercise}

\begin{snippetsolution}{linear-algebra-batch-6-ex-4-sol}{}
    \todo
\end{snippetsolution}

\begin{snippetexercise}{linear-algebra-batch-6-ex-5}{}
    In \(\mathbb A^3(\realnumbers)\) consider the lines \(l, l'\) with parametric equations
    \[
        \begin{cases}
            x = 3 + 2t \\
            y = -2 + t \\
            z = 4 + 5t
        \end{cases}, \quad \begin{cases}
            x = -3 + 2u \\
            y = 2-u \\
            z = 7-5u
        \end{cases}
    \]
    respectively.
    Determine whether the lines are coincident, parallel, intersecting or skew.
\end{snippetexercise}

\begin{snippetsolution}{linear-algebra-batch-6-ex-5-sol}{}
    \todo
\end{snippetsolution}

\end{document}